\subsection{The half-differential algebra question of scans 348--352}

\begin{theorem}[B601: \status{Refuted}; functorial exterior half-algebra]\label{sa08:B601}
Let $k\ne0$ be any commutative ring, $\Phi_n=k^{n+1}$, $T=1$,
$\Psi=\Phi/kT$, and $C^i=\Lambda^{i+1}\Phi$ with differential
$d=T\wedge-$. Identify its cocycles with $Z^i=\Lambda^i\Psi$ by
$z\mapsto T\wedge\widetilde z$. There is no natural simplicial
$k$-bilinear operation
\[
 \sigma_{0,1}:C^0\otimes_k Z^1\longrightarrow C^1
\]
that extends the exterior product on $Z^0\otimes Z^1$ and satisfies
$d\sigma(c,z)=(dc)\wedge z$. Consequently the proposed
half-differential algebra structure does not exist.
\end{theorem}
\begin{proof}
On $\Delta[1]$, write $c=(c_0,c_1)$ and represent $\bar v\in\Psi_1$
by $(v_0,v_1)$. Both $\Psi_1$ and $\Lambda^2\Phi_1$ have rank one,
so the most general bilinear operation has the form
\[
 \sigma(c,\bar v)_{01}=(a c_0+b c_1)(v_1-v_0),\qquad a,b\in k.
\]
The marked-unit condition gives $a+b=1$. Naturality on the three
edges of $\Delta[2]$ forces this formula on every edge. Thus
\begin{align*}
 (d\sigma(c,\bar v))_{012}
 &=a(c_1-c_0)(v_2-v_1)+b(c_1-c_2)(v_1-v_0),\\
 ((dc)\wedge\bar v)_{012}
 &=(c_1-c_0)(v_2-v_1)+(c_1-c_2)(v_1-v_0).
\end{align*}
Taking $c=(0,1,1)$, $v=(0,0,1)$ gives $a=1$.
Taking $c=(0,0,1)$, $v=(0,1,1)$ gives $b=1$.
Together with $a+b=1$ these imply $1=0$. The argument uses only
simplices of dimensions one and two and holds with arbitrary torsion.

The supplied program \texttt{code/C\_finite\_obstructions.py}
exhausts all coefficients and all pairs of functions on the
two-simplex over $\Z/2$, $\Z/3$, and $\Z/4$. It performs respectively
$256$, $6561$, and $65536$ Leibniz checks. In each ring the only
Leibniz-compatible pair is $(a,b)=(1,1)$, which violates the unit
condition. The universal certificate is the polynomial identity
\[
 (a+b-1)-(a-1)-(b-1)=1.\qedhere
\]
\end{proof}

\begin{proposition}[B602: \status{Proved}; corrected extension and ambiguity]\label{sa08:B602}
For a fixed pure exact sequence $0\to kT\to\Phi\to\Psi\to0$ with
$\Psi$ flat in a ringed topos, the individual lifting problem on
scan~352 has an obstruction in
\[
 \Ext^1_k(Z^{i+1}\otimes_k Z^j,Z^{i+j}).
\]
When that obstruction vanishes, its solutions form a torsor under
$\Hom_k(Z^{i+1}\otimes_k Z^j,Z^{i+j})$. Associativity imposes further
conditions. A specified splitting of $\Phi\to\Psi$ gives a
half-differential algebra, functorial under maps preserving that
splitting. In the standard simplicial example the solution set is
empty already at $(i,j)=(0,1)$.
\end{proposition}
\begin{proof}
The desired map is the middle map of a morphism between the exact
sequences
\[
\begin{aligned}
 0&\longrightarrow Z^i\otimes Z^j
 \longrightarrow C^i\otimes Z^j
 \longrightarrow Z^{i+1}\otimes Z^j\longrightarrow0,\\
 0&\longrightarrow Z^{i+j}\longrightarrow C^{i+j}
 \longrightarrow Z^{i+j+1}\longrightarrow0,
\end{aligned}
\]
whose outer maps are the prescribed exterior products. The first
sequence stays exact because $Z^j=\Lambda^j\Psi$ is flat. Push the
first extension out by its lower product and pull the second one
back by its upper product. A middle map exists precisely when these
extensions represent the same Yoneda class. This proves the
obstruction criterion. The difference of any two middle maps
vanishes on $Z^i\otimes Z^j$ and lands in $Z^{i+j}$, hence factors
uniquely through $Z^{i+1}\otimes Z^j$. Conversely adding any such
factor preserves the two outer maps. This proves the torsor claim.

If $s:\Psi\to\Phi$ is a supplied splitting, define
\[
 c\circ z=c\wedge\Lambda^j(s)(z),\qquad c\in C^i,\quad z\in Z^j,
\]
and define the left operations by graded commutativity in the
cochain degrees. Exterior associativity proves associativity
whenever two entries are cocycles; $T$ acts as the marked unit;
the restriction to cycles is the prescribed exterior product; and
$d=T\wedge-$ gives $d(c\circ z)=dc\circ z$. All constructions
commute with maps preserving $s$. The final assertion is
Statement~\ref{sa08:B601}.
\end{proof}

\begin{theorem}[B605: \status{Proved}; full ambiguity when the quotient has finite rank]\label{sa08:B605}
In Statement~\ref{sa08:B602}, suppose additionally that $\Psi$ is
locally free of constant finite rank $r$ in the ringed topos.
Then a half-differential algebra extending exterior multiplication
exists if and only if $0\to kT\to\Phi\to\Psi\to0$ splits.
If it exists, its full set of structures, including associativity,
is a torsor under
\[
 \prod_{i=0}^{r-1}
 \Hom_k(\Lambda^{i+1}\Psi,\Lambda^i\Psi).
\]
Equivalently, it is the set of independent retractions
$p_i:C^i\to Z^i$ of the cycle inclusions, for $0\le i<r$.
This additional finite-rank hypothesis does not hold for the
standard $\Psi$ as an object of the simplicial topos: its ranks
vary with the simplex dimension.
\end{theorem}
\begin{proof}
Write $A=\Lambda^*\Psi$ and $L=\Lambda^r\Psi$. The top term
$C^r=\Lambda^{r+1}\Phi$ is canonically $L$, by wedging with $T$.
Given a half-algebra, define $p_i(c)\in A^i$ by
\[
 p_i(c)\wedge z=c\circ z\in L,
 \qquad z\in A^{r-i}.
\]
The exterior pairing $A^i\otimes A^{r-i}\to L$ is perfect, so
this defines a unique sheaf morphism $p_i:C^i\to A^i$.
On cycles it is the identity. In particular $p_0(T)=1$, so
$p_0:\Phi\to k$ splits the marked line. For $i=r$, the map is
the fixed identity of $L$.

Associativity with two cycles implies
\[
 p_{i+j}(c\circ z)=p_i(c)\wedge z.
\]
Indeed pair both sides with an arbitrary element of $A^{r-i-j}$,
use $(c\circ z)\circ w=c\circ(z\wedge w)$, and use perfection
of the exterior pairing. The Leibniz condition also says that the
quotient $\overline d_i:C^i\to A^{i+1}$ is $A$-linear. Therefore
\[
 (p_i,\overline d_i):C^i\xrightarrow{\sim}
 A^i\oplus A^{i+1}
\]
identifies the entire right action with
$(a,b)\circ z=(a\wedge z,b\wedge z)$.

Conversely, choose arbitrary retractions $p_i$ of the cycle
inclusions, and use these displayed isomorphisms to define this
right action. The differential becomes $(a,b)\mapsto(b,0)$.
Define the left action by the required graded commutativity.
The exterior algebra identities now verify the unit, cycle product,
Leibniz rule, and associativity whenever two entries are cycles.
The top-degree pairing recovers exactly the chosen $p_i$, so this
construction and the preceding one are inverse.

If the original extension splits, its exterior powers give all
the required retractions, so the set is nonempty. The difference
of two retractions $C^i\to A^i$ factors uniquely through the
quotient $A^{i+1}$. Each choice set is consequently a torsor under
$\Hom_k(A^{i+1},A^i)$, and the choices are independent. This
proves both the existence criterion and the full stated ambiguity.
\end{proof}

The conditional ambiguity displayed on scan~352 is therefore
correct as a linear lifting statement, but existence must be checked
first. We now make the comparison with scans~357--361 precise.

\begin{theorem}[B603: \status{Known, with the stated consequences proved}; comparison with cup-one and Eilenberg--Zilber]\label{sa08:B603}
Let $C=N_*(X;k)$ be normalized simplicial chains. There are natural
equivariant maps
\[
 E(r)\otimes_k C\longrightarrow C^{\otimes r},\qquad
 E(r)=N_*(k[E\Sigma_r]),
\]
compatible with operadic composition. Each $E(r)$ is a free
$k[\Sigma_r]$-resolution of $k$. The degree-zero binary map is
Alexander--Whitney; its first homotopy is cup-one after dualizing.
These structures do not provide the half-algebra requested on
scan~352.
\end{theorem}
\begin{proof}
The existence and operadic compatibility of the displayed action
are Berger--Fresse, Theorem~2.1.1, with the construction in
\S2.2 of the opened arXiv source \cite{berger-fresse}.
We explain exactly what this theorem supplies here. In simplicial
degree $m$, the bar construction $E\Sigma_r$ consists of strings
$(\sigma_0,\ldots,\sigma_m)$. Faces delete an entry, degeneracies
repeat an entry, and the group acts by simultaneous translation.
That action is free on nondegenerate strings, so normalized chains
are free group modules. Inserting the identity at the beginning
contracts the augmented underlying chain complex; it need not be
equivariant to prove exactness. Thus $E(r)$ is the claimed free
resolution over every $k$.

For the binary generator $e_0=(\mathrm{id})$, the action is
\[
 x\longmapsto\sum_{a=0}^n
 x[0,\ldots,a]\otimes x[a,\ldots,n],\qquad x\in X_n.
\]
For $e_1=(\mathrm{id},\tau)$, we have
$\partial e_1=(\tau)-(\mathrm{id})$. If $h$ is the dual degree
$-1$ operation and $\mu$ the dual degree-zero product, the chain-map
identity is exactly
\[
 d h(a,b)+h(da,b)+(-1)^{|a|}h(a,db)
 =(-1)^{|a||b|}\mu(b,a)-\mu(a,b).
\]
This is a cup-one identity, with the sign convention specified by
this formula. Higher alternating strings give cup-$i$ operations,
as identified explicitly in Berger--Fresse, Theorem~2.1.1 and
the final paragraph of \S2.2. The free resolutions on scan~359
therefore have a precise operadic realization. Applying a cocycle
in each tensor factor and multiplying its values gives the usual
equivariant power construction indicated on scans~359 and~361.

By contrast, $\sigma_{ij}$ in Statement~\ref{sa08:B602} has
cohomological degree zero, is defined when one input is a cycle,
and must restrict to a specified exterior product. The cup-one
identity involves the difference of two products and has degree
minus one. Neither the Eilenberg--Zilber equivalence nor its higher
homotopies changes the three incompatible scalar equations in
Statement~\ref{sa08:B601}. This proves the claimed distinction.
\end{proof}

\begin{proposition}[B604: \status{Refuted}; pointwise reading of scan~360]\label{sa08:B604}
The complexes of chain maps appearing on scan~360 are not, for
each fixed $X$, free $\Sigma_r$-resolutions of $\Z$. The inserted
qualification ``for variable $X$'' is essential. A precise true
replacement is the natural operation complex together with the
free Barratt--Eccles resolution and its action in
Statement~\ref{sa08:B603}.
\end{proposition}
\begin{proof}
For $X=\Delta[0]$ and $r=2$, the normalized version of either
pointwise mapping complex is $\Z$ in degree zero with trivial
$\Sigma_2$-action, hence is not free over $\Z[\Sigma_2]$.
The unnormalized version fails freeness as well: its degree-zero
Hom module has an equivariant direct summand $\Z$, obtained by
retaining only the map on chain degree zero. The trivial module
$\Z$ is not projective over $\Z[\Sigma_2]$, since an equivariant
section of the augmentation would require $2a=1$ in $\Z$.

Acyclicity also fails. For
$X=\Delta[1]/\partial\Delta[1]$, let $C=N_*(X;\Z)
=\Z[0]\oplus\Z[1]$ with zero differential. Eilenberg--Zilber
identifies $N_*(X\times X)$ up to chain homotopy with
$T=C^{\otimes2}$, whose ranks are $(1,2,1)$ and whose differential
is zero. Either pointwise comparison Hom complex is therefore
chain-homotopy equivalent to $\underline\Hom(T,T)$.
Its degree-one homology has rank
$1\cdot2+2\cdot1=4$. Normalized and unnormalized chains are
chain-homotopy equivalent, so the conclusion holds for both
conventions on scan~357. The finite ranks are checked in
\texttt{code/C\_finite\_obstructions.py}. Berger--Fresse define
the Eilenberg--Zilber operad using \emph{natural transformations in
$X$}, not one pointwise mapping complex; see the paragraph
``The Eilenberg--Zilber operad'' after their Theorem~2.1.1.
\end{proof}
